\subsection{Retrieval (Phase A)}
\label{sec:retrieval}

\subsubsection{BM25 Retrieval: From PyTerrier PISA to pg\_textsearch}

In previous editions, we relied on the PISA indexing framework~\cite{mallia2019pisa} accessed through the PyTerrier interface for BM25-based first-stage retrieval. While effective, this approach presented several operational limitations. The PISA indexes required loading into memory at each pipeline iteration, incurring substantial disk I/O time. Additionally, the indexes existed as separate artifacts from the document storage, creating synchronization challenges.

For the 14th edition, we migrated BM25 retrieval to PostgreSQL's built-in pg\_textsearch module. This change integrated document storage and full-text search into a single source of truth, eliminating the need for separate index files. The PostgreSQL GiN (Generalized Inverted Index) indexes are maintained automatically alongside document updates, and query execution benefits from decades of query planner optimization. Retrieval speed improved notably, particularly for batch queries, as the database engine handles concurrent access efficiently without the memory-loading overhead of PISA.

\subsubsection{Dense Retrieval with Qdrant}

A major architectural shift was the adoption of Qdrant as the vector database for dense embedding indexing. In prior editions, we precomputed and stored all cross-embedding similarities as flat NumPy arrays, a practice that became increasingly untenable as the document corpus grew. This approach consumed substantial disk space and offered no flexibility for ad-hoc similarity searches or incremental updates.

Qdrant provides several advantages over the previous approach. First, GPU acceleration for similarity search enables significantly faster retrieval compared to CPU-based cosine similarity computation on NumPy arrays. Second, the HNSW (Hierarchical Navigable Small World) indexing in Qdrant provides an efficient approximate nearest neighbor search with configurable precision-speed tradeoffs. Third, storage efficiency improves markedly, as Qdrant compactly stores embeddings and indexes rather than materializing the full similarity matrix. Finally, the Qdrant API supports dynamic collection management, enabling filtering by metadata fields and seamless updates as new documents are indexed.

\subsubsection{Text Embeddings Interface (TEI)}

For embedding generation, we adopted the Text Embeddings Interface (TEI) from Hugging Face, a fast, Rust-based implementation optimized for serving text embeddings. TEI provides tokenization and model inference in a single, efficient binary, supporting both CPU and GPU execution. Compared to loading models directly through the Transformers library in Python, TEI reduced embedding generation latency substantially, enabling us to index the full PubMed corpus more efficiently.

\subsubsection{Query Expansion: HyDE and Context-1}

We explored two query expansion strategies to improve retrieval relevance. The first, Hypothetical Document Embeddings (HyDE), replaces the original query with a hypothetical answer generated by an LLM. This hypothetical answer is then used as the query for dense retrieval, leveraging the semantic similarity between the generated answer and genuinely relevant documents. The intuition is that an LLM-generated answer, even if factually incomplete, shares the linguistic and conceptual structure of relevant documents more closely than the original question.

The second strategy, Context-1~\cite{context1}, employs an LLM within the standard RAG pipeline to iteratively refine the retrieval process. The LLM analyzes the retrieved context and generates a refined representation of the information need, which is then used to retrieve additional documents. This feedback loop allows the system to dynamically adapt its retrieval strategy based on the content of already-retrieved documents.

\subsubsection{Reranker Training Pipeline}

We developed a new, flexible training pipeline for cross-encoder rerankers. The pipeline supports multiple base models, training strategies, and data configurations. A notable improvement was the incorporation of dense retrieval for negative sampling. In previous editions, negative examples for reranker training were selected exclusively from BM25 results, which limited the diversity of negatives. By including dense retrieval results in the negative generation process, the reranker is exposed to a broader range of challenging distractors, improving its ability to distinguish relevant from irrelevant documents across different retrieval paradigms.

We trained a total of 29 reranker models based on diverse architectures including BGE-reranker-v2-m3, PubMedBERT, BioLinkBERT, BioBERT, MedCPT, and the LLaMA Nemotron reranker, among others. Training employed a pairwise loss with varying configurations of epochs, samplers, and hard negative counts.

\subsubsection{Fusion Strategies}

For combining BM25 and dense retrieval scores, we evaluated both Reciprocal Rank Fusion (RRF) and weighted sum approaches. RRF, which combines rankings based solely on reciprocal positions rather than raw scores, received particular attention due to its robustness to score scale differences between heterogeneous retrieval methods. The RRF score for a document $d$ is computed as:

\begin{equation}
RRF(d) = \sum_{r \in R} \frac{1}{k + rank_r(d)}
\end{equation}

where $R$ is the set of retrieval runs and $k$ is a constant (typically 60) that mitigates the impact of high rankings on the final score.

\subsubsection{Snippet Generation}

For the first time, we participated in the snippet generation subtask. Given a question and a set of retrieved documents, the system must identify the most relevant text spans within those documents. We approached this as a fine-grained retrieval problem, using the same reranker models trained for document retrieval to score individual sentences or passages within the top-ranked documents. The highest-scoring passages were selected as snippets.

\subsubsection{Future Directions: SPLARe and ColBERT}

As part of our ongoing research, we are exploring two additional approaches not yet deployed in competition. SPLARe (Sparse Late Interaction for Retrieval)~\cite{formal2021splade} combines the efficiency of sparse representations with the expressiveness of learned term weighting, offering a middle ground between lexical and dense retrieval. ColBERT~\cite{khattab2020colbert} employs a late interaction paradigm where query and document tokens interact only at the final scoring stage, enabling efficient indexing while preserving fine-grained semantic matching. Both approaches represent promising directions for future editions.
